\pagestyle{plain}

This section provides an overview of a few directions we are interested in investigating in the near future. The reader is invited to have a look at these various or related problems.

\subsection{Cokernel-Kernel operators}

This section introduces a closure operator on subcategories of $(\mathscr{A},\mathcal{E})$ that is slightly different from the Gen-Sub operator. More precisely, given a subcategory $\mathscr{C}$, we will construct the smallest subcategory $\mathscr{D} \supseteq \mathscr{C}$ such that:
\begin{enumerate}[label=$\bullet$,itemsep=1mm]
    \item $\mathscr{D}$ is closed under cokernels of $\mathcal{E}$-monomorphisms in $\mathscr{D}$; and,
    \item  $\mathscr{D}$ is closed under kernels of $\mathcal{E}$-epimorphisms in $\mathscr{D}$.
\end{enumerate}

The \new{$\Coker_\mathcal{E}$-operator} on subcategories of $\mathscr{A}$ is defined as follows: given a subcategory $\mathscr{C} \subseteq \mathscr{A}$, we define $\Coker_\mathcal{E}(\mathscr{C})$ as the full subcategory, closed under sums and summands, of $\mathscr{A}$ containing the cokernel of any $\mathcal{E}$-monomorphism $f : X \longrightarrow Y$ such that $X,Y \in \mathscr{C}$. The \new{$\Ker_\mathcal{E}$-operator} on subcategories of $\mathscr{A}$ is defined analogously: given a subcategory $\mathscr{C} \subseteq \mathscr{A}$, we define $\Ker_\mathcal{E}(\mathscr{C})$ as the full subcategory, closed by sums and summands, of $\mathscr{A}$ containing the kernel of any $\mathcal{E}$-epimorphism $g : Y \longrightarrow X$  such that $Y,X \in \mathscr{C}$. Given an object $M \in \mathscr{C}$, by abusing of notation, we set $\Coker_\mathcal{E}(M) = \Coker_\mathcal{E}(\add(M))$ and $\Ker_\mathcal{E}(M) = \Ker_\mathcal{E}(\add(M))$.

Given a subcategory $\mathscr{C} \subseteq \mathscr{A}$, we define a sequence of subcategories $(\CK_{\mathcal{E}}^i(\mathscr{C}))_{i\in \mathbb{N}}$ as it follows:
\begin{enumerate}[label=$\bullet$, itemsep=1mm]
    \item we set $\CK_{\mathcal{E}}^0(\mathscr{C}) = \mathscr{C}$; and,
    \item for all $i \geqslant 1$, we define $\CK^{i}_{\mathcal{E}}(\mathscr{C})$ to be the subcategory of $\mathscr{A}$ additively generated by objects in both $\Coker_{\mathcal{E}}(\CK^{i-1}_{\mathcal{E}}(\mathscr{C}))$ and $\Ker_{\mathcal{E}}(\CK^{i-1}_{\mathcal{E}}(\mathscr{C}))$.
\end{enumerate}

By the previous remark, $(\CK_{\mathcal{E}}^i(\mathscr{C}))_{i\in \mathbb{N}}$ is an increasing sequence of additive subcategories of $\mathscr{A}$. Therefore, as $\mathscr{A}$ is abelian, this sequence of subcategories admits a colimit.

\begin{definition} \label{def:CKop}
    The \new{$\CK_\mathcal{E}$-operator} on subcategories of $\mathscr{A}$ is defined as follows: given a subcategory $\mathscr{C} \subseteq \mathscr{A}$, we define $\CK_{\mathcal{E}}(\mathscr{C})$ as the colimit of the sequence of subcategories $(\CK_\mathcal{E}^i(\mathscr{C}))_{i\in \mathbb{N}}$. As before, we set $\CK_\mathcal{E}(M) = \CK_\mathcal{E}(\add(M))$ by abuse of notations.
\end{definition}

\begin{remark}
It is clear that $\CK_{\mathcal{E}}(\mathscr{C}) \subseteq \GS_{\mathcal{E}}(\mathscr{C})$. Therefore, by~\cref{prop:EAdapt}, $\CK_{\mathcal{E}}(T)$ is $\mathcal{E}$-adapted for any rigid object $T$.

\end{remark}

\begin{ex} \label{ex:A7typeCK} In \cref{fig:CKCalc1}, we calculate $\CK_{\mathcal{E}_\diamond}(\mathscr{C})$ for a subcategory $\mathscr{C}$ of $\rep(Q)$ and a fixed quiver $Q$ of $A_7$ type. This example illustrates the slight difference between $\CK_{\mathcal{E}_\diamond}(\mathscr{C})$ and $\GS_{\mathcal{E}_\diamond}(\mathscr{C})$.

\begin{figure}[!ht]
    \leavevmode\\[-0.5ex]  % forces heading to appear before content
    \centering
    \begin{tikzpicture}
				\begin{scope}[line width=.5mm,->, >= angle 60]
                    \node at (-.6,0){$Q=$};
					\node (a) at (0,0){$1$};
					\node (b) at (1,0){$2$};
					\node (c) at (2,0){$3$};
					\node (d) at (3,0){$4$};
                    \node (e) at (4,0){$5$};
                    \node (f) at (5,0){$6$};
                    \node (g) at (6,0){$7$};
					\draw (a) -- (b);
					\draw (c) -- (b);
					\draw (c) -- (d);
                    \draw (d) -- (e);
                    \draw (f) -- (e);
                    \draw (f) -- (g);
				\end{scope}
                \begin{scope}[yshift=-2cm, xshift=-.5cm, line width=.3mm, ->, >= angle 60]
					\node (2) at (0,0){\scalebox{.7}{$\llrr{2}$}};
					\node (12) at (1,1){\scalebox{.7}{$\llrr{1,2}$}};
					\node (2345) at (1,-1){\scalebox{.7}{$\llrr{2,5}$}};
					\node (45) at (0,-2){\scalebox{.7}{$\llrr{4,5}$}};
                    \node[circle,line width=0.2mm,double,fill=lava!20,draw] (5) at (-1,-3){\scalebox{.7}{$\llrr{5}$}};
                    \node[circle,line width=0.2mm,double,fill=lava!20,draw] (567) at (0,-4){\scalebox{.7}{$\llrr{5,7}$}};
                    \node (7) at (-1,-5){\scalebox{.7}{
							$\llrr{7}$}};
                    \node (56)[circle,line width=0.2mm,double,fill=lava!20,draw] at (1,-5){\scalebox{.7}{$\llrr{5,6}$}
                                };
                     \node (4) at (3,-5){\scalebox{.7}{$\llrr{4}$}};
                     \node[circle,line width=0.2mm,double,fill=lava!20,draw] (23) at (5,-5){\scalebox{.7}{$\llrr{2,3}$}};
                    \node (1) at (7,-5){\scalebox{.7}{$\llrr{1}$}};
                    \node (456) at (2,-4){\scalebox{.7}{$\llrr{4,6}$}};
                    \node (234) at (4,-4){\scalebox{.7}{$\llrr{2,4}$}};
                     \node (123) at (6,-4){\scalebox{.7}{$\llrr{1,3}$}};
                     \node (4567) at (1,-3){\scalebox{.7}{$\llrr{4;7}$}};
                     \node(23456) at (3,-3){\scalebox{.7}{$\llrr{2,6}$}};
                     \node (1234) at (5,-3){\scalebox{.7}{$\llrr{1,4}$}};
                     \node (3) at (7,-3){\scalebox{.7}{$\llrr{3}$}};
					\node[circle,line width=0.2mm, double,fill=lava!20,draw] (12345) at (2,0){\scalebox{.7}{$\llrr{1,5}$}};
					\node (234567) at (2,-2){\scalebox{.7}{$\llrr{2,7}$}};
                    \node[circle,line width=0.2mm,fill=darkgreen!20,draw] (123456) at (4,-2){\scalebox{.7}{$\llrr{1,6}$}};
					\node (34) at (6,-2){\scalebox{.7}{$\llrr{3,4}$}};
					\node[circle,line width=0.2mm,fill=darkgreen!20,draw] (1234567) at (3,-1){\scalebox{.7}{$\llrr{1,7}$}};
                    \node[circle,line width=0.2mm,fill=darkgreen!20,draw] (3456) at (5,-1){\scalebox{.7}{$\llrr{3,6}$}};
                    \node[circle,line width=0.2mm,fill=darkgreen!20,draw] (34567) at (4,0){\scalebox{.7}{$\llrr{3,7}$}};
                    \node (6) at (6,0){\scalebox{.7}{$\llrr{6}$}};
                    \node[circle,line width=0.2mm,double,fill=lava!20,draw] (345) at (3,1){\scalebox{.7}{$\llrr{3,5}$}};
                    \node (67) at (5,1){\scalebox{.7}{$\llrr{6,7}$}};
					\draw (2) -- (12);
					\draw (2) -- (2345);
					\draw (45) -- (2345);
					\draw (12) -- (12345);
					\draw (2345) -- (12345);
					\draw (2345) -- (234567);
                    \draw (12345) -- (345);
                    
                    
                    \draw (5) -- (45);
                    \draw (7) -- (567);
                    \draw (567) -- (4567);
                    \draw (4567) -- (234567);
                    \draw (234567) -- (1234567);
                    \draw (1234567) -- (34567);
                    \draw (34567) -- (67);
                    \draw (56) -- (456);
                    \draw (456) -- (23456);
                    \draw (23456) -- (123456);
                    \draw (123456) -- (3456);
                    \draw (3456) -- (6);
                    \draw (4) -- (234);
                    \draw (234) -- (1234);
                    \draw (1234) -- (34);
                    \draw (23) -- (123);
                    \draw (123) -- (3);
					
					\draw (5) -- (567);
                    \draw (567) -- (56);
                    \draw (45) -- (4567);
                    \draw (4567) -- (456);
                    \draw (456) -- (4);
                    \draw (234567) -- (23456);
                    \draw (23456) -- (234);
                    \draw (234) -- (23);
                    \draw (12345) -- (1234567);
                    \draw (1234567) -- (123456);
                    \draw (123456) -- (1234);
                    \draw (1234) -- (123);
                    \draw (123) -- (1);
                    \draw (345) -- (34567);
                    \draw (34567) -- (3456);
                    \draw (3456) -- (34);
                    \draw (34) -- (3);
                    \draw (67) -- (6);	
				\end{scope}
    \end{tikzpicture}
    \caption[fragile]{Calculation of $\CK_{\mathcal{E}_\diamond}(\mathscr{C})$ where $\mathscr{C} = \add \left( \protect\begin{tikzpicture}[baseline={(0,-.1)}]
        \node[circle, line width=0.2mm, double, fill=lava!20,minimum size=1em,draw] at (0,0){$ $};
    \end{tikzpicture}\right)$. We get $\CK_{\mathcal{E}_\diamond}^1(\mathscr{C}) = \add \left(\protect\begin{tikzpicture}[baseline={(0,-.1)}]
        \node[circle, line width=0.2mm, double, fill=lava!20,minimum size=1em,draw] at (0,0){$ $};
    \end{tikzpicture} + \protect\begin{tikzpicture}[baseline={(0,-.1)}]
        \node[circle, line width=0.2mm, fill=darkgreen!20,minimum size=1em,draw] at (0,0){$ $};
    \end{tikzpicture}\right)$. Here, we obtain that $\CK_{\mathcal{E}_\diamond}(\mathscr{C}) = \CK_{\mathcal{E}_\diamond}^1(\mathscr{C})$.}
    \label{fig:CKCalc1}
\end{figure}
\end{ex}

For the same subcategory $\mathscr{C}$, $\GS_{\mathcal{E}_\diamond}(\mathscr{C})$ still corresponds to the subcategory shown in Figure~\ref{fig:GSCalc1}. By adding the object $\llrr{1,3}$ to $\mathscr{C}$, thereby making it maximally rigid, we obtain $\GS_{\mathcal{E}_\diamond}(\mathscr{C}) = \CK_{\mathcal{E}_\diamond}(\mathscr{C})$. We observe that, in general, a minimal basic object $T$ for which $\GS_{\mathcal{E}}(T)$ is a maximally $\mathcal{E}$-adapted extension-closed subcategory does not necessarily have to be tilting. This suggests that studying the $\CK_\mathcal{E}$-operator could be valuable.

\begin{remark}
By minimal, we mean that no direct summand can be removed from $T$ basic without losing the property that $\GS_{\mathcal{E}}(T)$ is a maximally $\mathcal{E}$-adapted extension-closed subcategory.
\end{remark}

\begin{definition}
We recall that a subcategory $\mathscr{D}$ is called \new{$\mathcal{E}$-wide} if, for any short $\mathcal{E}$-exact sequence
\[
\begin{tikzcd}
	\xi:\quad 0 \arrow[r] & E \arrow[r, tail] & F \arrow[r, two heads] & G \arrow[r] & 0,
\end{tikzcd}
\]
whenever two of the three objects $E$, $F$, or $G$ belong to $\mathscr{D}$, then so does the third.
\end{definition}

\begin{prop}
Let $(\mathscr{A}, \mathcal{E})$ be an exact category with $\mathscr{A}$ hereditary. Then the following statements hold:
\begin{enumerate}[label=$(\alph*)$, itemsep=1mm]
    \item \label{aCKGS} if $T \in \mathscr{A}$ is rigid, then $\CK_{\mathcal{E}}(T)$ is an $\mathcal{E}$-adapted and $\mathcal{E}$-wide subcategory;
    \item \label{bCKGS} if $T \in \mathscr{A}$ is tilting, then $\CK_{\mathcal{E}}(T) = \GS_{\mathcal{E}}(T)$. 
\end{enumerate}
\end{prop}

\begin{proof}
The proof of \ref{aCKGS} follows as in \cref{prop:Extclosed}, noting that the subcategory is no longer necessarily closed under arbitrary $\mathcal{E}$-subobjects and $\mathcal{E}$-quotients, but it remains closed under taking kernels of $\mathcal{E}$-epimorphisms and cokernels of $\mathcal{E}$-monomorphisms within itself.

The proof of \cref{thm:maximal} can be adapted to the setting of $\CK_{\mathcal{E}}(T)$ for showing \ref{bCKGS}. Together with the result from \ref{aCKGS}, this implies that both $\CK_{\mathcal{E}}(T)$ and $\GS_{\mathcal{E}}(T)$  are maximal $\mathcal{E}$-adapted extension-closed subcategories with $\CK_{\mathcal{E}}(T) \subseteq \GS_{\mathcal{E}}(T)$, which yields the desired equality.
\end{proof}

\begin{conj} \label{conj:TiltasmingenCK}
Let $(\mathscr{A},\mathcal{E})$ be an exact category with $\mathscr{A}$ hereditary. Then, $T \in \mathscr{A}$ is minimal such that $\CK_{\mathcal{E}}(T)$ forms a maximally $\mathcal{E}$-adapted extension-closed subcategory if and only if $T$ is tilting.  
\end{conj}

\subsection{Congruences on the tilting lattice} \label{ss:tiltposet} The equivalence relation $\approx_{\mathcal{E}}$ seems to hide other interesting properties in the lattice $(\Tilt(Q), \Tleq)$. Considering the equivalence relation $\approx_{\mathcal{E}}$, we can ask ourselves if it defines well-behaved lattice quotients, usually called \emph{lattice congruences}. 
\begin{definition} \label{def:latticecongr}
    Let $(L, \leqslant, \vee, \wedge)$. A \new{congruence} of $L$ is an equivalence relation $\sim$ on $L$ such that, for any $a,b,x,y \in L$, if $a \sim x$ and $b \sim y$, then both $a \vee b \sim x \vee y$ and $a \wedge b \sim x \wedge y$. The \new{quotient lattice} of $L$ on $\sim$ is the lattice $(L / \sim, \leqslant, \vee, \wedge)$ induced by $L$.
\end{definition}
N. Reading gives a criterion on equivalence relations that are congruences on a lattice without using join and meet operations \cite{R04}.

For any $n \in \mathbb{N}$, denote by $\mathcal{B}_n$ the Boolean lattice of order $n$. In \cref{fig:posettiltEdiamond}, we can check that $\approx_{\mathcal{E}_\diamond}$ defines a congruence on $(\Tilt(Q), \Tleq)$, and the lattice congruence obtained is isomorphic to the boolean lattice $\mathcal{B}_3$. In light of this example, we made a stronger conjecture that encapsulates the general behavior of $\approx_{\mathcal{E}_\diamond}$.

\begin{conj} \label{conj:latticequotient} Let $Q$ be an $A_n$ type quiver for some $n \in \mathbb{N}^*$. The equivalence relation $\approx_{\mathcal{E}_\diamond}$ defines a congruence on $(\Tilt(Q), \Tleq)$, and the obtained lattice congruence quotient is isomorphic to the Boolean lattice $\mathcal{B}_{n-1}$  
\end{conj}

\begin{example} \label{ex:latticequotientEdiamond}
  In \cref{fig:posettiltEdiamond}, the equivalence relation $\approx_{\mathcal{E}_{\diamond}}$ gives a lattice congruence, given by identifying gathered tilting representations, and the quotient is isomorphic to $\mathcal{B}_3$.  
\end{example}


A rough plan to prove this conjecture is as follows. First, we show that $\approx_{\mathcal{E}_\diamond}$ defines a congruence on $(\Tilt(Q), \Tleq)$ by using Reading's criterion. Then, we can easily count the number of equivalence classes of $\approx_{\mathcal{E}_\diamond}$ in $\Tilt(Q)$.

\begin{lemma} \label{lem:numberofeqclassEdiamond}
    Let $Q$ be an $A_n$ type quiver for some $n \in \mathbb{N}^*$. There are exactly $2^{n-1}$ equivalence classes for $\approx_{\mathcal{E}_\diamond}$ in $\Tilt(Q)$.
\end{lemma}

\begin{proof}
    By \cref{thm:CJRmaxcomb,th:mCJR=GS}, those equivalence classes are parametrized by pairs $(\B,\E)$ such that $\{\B,\E[1]\}$ is a bipartition of $\{1,\ldots,n+1\}$. Then there is a one-to-one correspondence $\phi$ from subsets of $\{2,\ldots,n\}$ to such pairs given by \[\forall \A \subseteq \{2,\ldots,n\},\ \phi(\A) = (\{1\} \cup \A,   \{2,\ldots, n+1\}) \setminus \A.\]
    We got the desired result.
\end{proof}

Finally, we can construct the lattice isomorphism explicitly. Note that the order of subsets of $\{2, \ldots, n\}$ induced by $(\Tilt(Q) / \approx_{\mathcal{E}_\diamond}, \Tleq)$ does not coincide with the usual inclusion order on those subsets. 

\begin{remark} \label{rem:trueforlinear}
   If $Q$ is a linearly-oriented $A_n$ type quiver, then $(\Tilt(Q), \Tleq)$ is isomorphic to the classical order on \emph{binary trees} with $n$ internal nodes, and $\approx_{\mathcal{E}_\diamond}$ identifies trees by their \emph{canopy} \cite{V09}. In this special case, the conjecture was proved by J.-L. Loday and M. O. Racia  \cite{LR98,L04}. 
\end{remark}

We aim to prove this conjecture soon. We can also try to characterize exact structures $\mathcal{E}$ on a type $ADE$ quiver $Q$ such that $\approx_\mathcal{E}$ defines a congruence on $\Tilt(Q)$. It could provide a deeper understanding of a subfamily of well-known lattice congruences within the Cambrian lattice.


